While studying certain cohomological properties of Koszul rings in~\cite{jps}, the authors were led to a new notion, the so-called \emph{Koszul pairs}, which we explain in brief. Let $R$ be a semisimple ring. A graded $R$-ring is a graded algebra in the tensor category of $R$-bimodules with respect to the tensor product (of bimodules). A graded $R$-ring $A = \op_{n \in \bb{N}} A^n$ is called \emph{connected} if $A^0 = R$. Connected graded $R$-corings are def\/ined by duality. By def\/inition, an \emph{almost Koszul pair} consists of a graded connected $R$-ring $A$ and a graded connected $R$-coring $C$, together with an $R$-bimodule isomorphism $\theta_{A,C}\colon C_1 \to A^1$. These data must be compatible, in the sense that the composition of the three maps below must be zero:
\begin{gather}\label{ec:theta}
C_2 \xrar{\Delta} C_1 \ot C_1 \xrar{\theta_{C,A} \ot \theta_{C,A}} A^1 \ot A^1 \xrar{\mu^{1,1}} A^2,	
\end{gather}
where $\Delta_{1,1}$ and $\mu^{1,1}$ denote the components of (co)multiplication maps for $C$ and $A$, respectively.
